Two natural comparisons help explain both the model and the limits of direct
representations. As temporal correlation approaches the static case, one can
compare moment coefficients and their growth rates. One may also try to copy a
Brownian clock representation from the first moment to the second. The first
comparison yields useful inequalities; a low-order calculation obstructs the
particular clock construction considered here.

Recall the mixed expansion from Chapter~\ref{S:G0}, the upper bound from
Chapter~\ref{S:book-independent-upper}, and the source-version qualifications in
Section~\ref{S:known}. Section~\ref{S:near-static-lower} compares the canonical
solution with its static counterpart and records the near-static mean
comparison. Temporal interpolation then gives log-convexity of the upper growth
rates under the stated assumptions. Section~\ref{S:two-clock-screen} tests
direct Brownian clock copying at small order and quantifies the discrepancy in
the cross-pair transform. This obstruction concerns that proposed
representation, not all possible two-clock constructions. Neither the static
comparison nor the clock test identifies the intrinsic constant proved in
Chapter~\ref{S:book-constant}.
